\documentclass[journal]{IEEEtran}
\usepackage{cite}
\usepackage{amsmath,amssymb,amsfonts}
\usepackage{graphicx}
\usepackage{booktabs}
\usepackage{array}
\usepackage{url}
\usepackage{xcolor}

\begin{document}

\title{ATRNet-STAR: Depression Angle Diversity, Few-Shot Recognition,
and Cross-Dataset Generalization in SAR Automatic Target Recognition}

\author{Mousa~Shaya,~\IEEEmembership{}
\thanks{M. Shaya is with Shaya Software and Systems LLC, Acton, MA, and
Arizona State University, Tempe, AZ (e-mail: mshaya01.ml@gmail.com).}}

\maketitle

\begin{abstract}
Synthetic Aperture Radar (SAR) Automatic Target Recognition (ATR) requires
models that generalize across depression angles, sensor configurations, and
datasets, yet labeled SAR imagery is scarce and expensive to acquire. We
present a systematic study of these generalization challenges using the
ATRNet-STAR dataset (40 vehicle classes, 4 depression angles) and the MSTAR
benchmark (10 military vehicle classes). We make six contributions:
(1)~a depression angle ablation showing that adding a single 60\textdegree{}
view to a 15\textdegree{}-only baseline improves EOC accuracy by +30.4~pp;
(2)~a data efficiency analysis showing that multi-angle training with as few
as 25 images per vehicle per angle surpasses a model trained on thousands of
single-angle images; (3)~a finding that neither angle-weighted sampling nor
test-time augmentation (TTA) improves EOC performance---TTA reduces accuracy
by 24.3~pp---confirming that angle invariance requires real multi-angle data;
(4)~a few-shot prototypical classifier achieving 98\%+ accuracy with only 5
reference images per class, replicated on ATRNet-STAR and MSTAR; (5)~a
cross-dataset transfer experiment revealing a 38~pp SAR-to-SAR domain gap
(61.1\% ATRNet-STAR$\to$MSTAR vs.\ 99.1\% MSTAR-specific); and (6)~a
demonstration that optical pretrained features (DINOv2, 6.1\%) are
ineffective for SAR ATR, motivating domain-specific approaches.
\end{abstract}

\begin{IEEEkeywords}
synthetic aperture radar, automatic target recognition, depression angle,
few-shot learning, prototypical networks, domain adaptation, ConvNeXt
\end{IEEEkeywords}

% ── I. Introduction ────────────────────────────────────────────────────────────
\section{Introduction}

SAR ATR is operationally critical: SAR sensors image through clouds and at
night, making them uniquely valuable for persistent surveillance. Deploying
ATR systems in practice surfaces three generalization challenges that
laboratory benchmarks often obscure.

\textit{The depression angle problem.} A vehicle imaged at 15\textdegree{}
looks fundamentally different from the same vehicle at 60\textdegree{} due to
layover, foreshortening, and shadow geometry. Most published results report
Standard Operating Condition (SOC) accuracy---train and test at the same
angle---which inflates operational expectations. Extended Operating Condition
(EOC) evaluation, where the test angle is withheld, is a harder and more
realistic measure.

\textit{The data scarcity problem.} Unlike optical vision, SAR ATR datasets
are measured in thousands of examples. This motivates few-shot approaches.

\textit{The domain gap problem.} Features learned on one SAR dataset do not
transfer freely to another. Understanding the magnitude of these gaps is
essential for building systems that generalize.

We address all three with a unified experimental framework. Our primary
dataset is ATRNet-STAR~\cite{atrnet}, a large-scale SAR vehicle dataset
collected at four depression angles (15\textdegree, 30\textdegree,
45\textdegree, 60\textdegree) across 40 vehicle classes. We supplement with
MSTAR~\cite{mstar}, the canonical 10-class military vehicle benchmark. Our
backbone is ConvNeXt-Small~\cite{convnext}, fine-tuned from ImageNet weights,
which substantially outperforms DINOv2~\cite{dinov2} on SAR imagery.

% ── II. Related Work ───────────────────────────────────────────────────────────
\section{Related Work}

\textbf{SAR ATR benchmarks.} MSTAR~\cite{mstar} remains the most widely
cited benchmark, with CNN accuracies of 97--99\% on the standard SOC task and
a well-studied EOC protocol. ATRNet-STAR~\cite{atrnet} extends this to 40
classes across four depression angles, spanning up to 45\textdegree{} of
angular change versus the 13\textdegree{} change in MSTAR's standard EOC.

\textbf{Depression angle generalization.} Performance degradation under
angular variation is well documented~\cite{atrnet}. Sparse representation,
complex-valued networks, and multi-view transformers have been applied, but
systematic study of which angles contribute most and how many samples per
angle are needed has not been reported.

\textbf{Few-shot SAR ATR.} Prototypical networks~\cite{proto} have been
applied to SAR ATR. DGP-Net~\cite{dgpnet} proposes a dense graph convolutional
network for few-shot MSTAR tasks, improving prior methods by over 4\%. We
differ in operating on ATRNet-STAR's 40-class problem and show that a frozen
backbone with nearest-centroid classification achieves 98\%+ accuracy with 5
reference images---no specialized architecture required.

\textbf{Foundation models for SAR.} SARATR-X~\cite{sarx} trains a
self-supervised backbone on 0.18M unlabeled SAR samples. Our DINOv2
experiment provides complementary motivation: optical SSL pretraining (6.1\%)
is essentially useless for SAR, confirming that domain-specific approaches are
necessary.

% ── III. Datasets and Methods ──────────────────────────────────────────────────
\section{Datasets and Methods}

\subsection{ATRNet-STAR}
ATRNet-STAR contains SAR chips of 40 ground vehicle classes at four depression
angles (15\textdegree, 30\textdegree, 45\textdegree, 60\textdegree). Images
are single-polarization magnitude chips converted to 8-bit grayscale via
log-magnitude normalization. \textit{SOC:} 70/15/15 train/val/test split
stratified by class. \textit{EOC:} train on 15\textdegree{} only, test on
\{30\textdegree, 45\textdegree, 60\textdegree\}.

\subsection{MSTAR}
MSTAR contains 10 military vehicle classes at two angles: 17\textdegree{}
(training) and 15\textdegree{} (test), stored as Phoenix binary complex
float32. Standard SOC split: 3,128 train / 543 val / 3,203 test.

\subsection{Model Architecture}
We use ConvNeXt-Small pretrained on ImageNet-1K. The classification head is
replaced with a linear layer sized to the target number of classes. Training
uses AdamW (lr=$10^{-4}$) with cross-entropy loss and random
horizontal/vertical flips and rotation augmentation.

\subsection{Prototypical Classification}
We extract features from the penultimate layer using a forward hook. Given $N$
reference images per class, per-class prototypes are computed by averaging
feature vectors. Test samples are classified by nearest centroid under
Euclidean distance. The backbone is frozen---no test-time training is required.

\subsection{Code Availability}
All code is publicly available at \url{https://github.com/mousashaya/sar-atr-ship}.

% ── IV. Experiments ────────────────────────────────────────────────────────────
\section{Experiments}

\subsection{SOC Baselines}
\textbf{ATRNet-STAR SOC} (40 classes, 4 angles): 97.2\% test accuracy.
\textbf{MSTAR SOC} (10 classes, 17\textdegree{}$\to$15\textdegree{}): 99.1\%
test accuracy, matching the published 97--99\% range~\cite{mstar}.

\subsection{EOC Baseline and Angle Ablation}

Fig.~\ref{fig:umap} visualizes the feature space under SOC and EOC backbones.
The SOC backbone places test images near their correct vehicle clusters;
the EOC backbone (trained on 15\textdegree{} only) produces an undifferentiated
scatter at 43.7\% accuracy---a 53.5~pp gap versus the SOC ceiling.

\begin{figure}[t]
  \centering
  \includegraphics[width=\columnwidth]{soc_eoc_overlay.png}
  \caption{UMAP of the 768-D feature space. Gray: 15\textdegree{} reference
  images. Colored: test images at 30\textdegree, 45\textdegree, 60\textdegree.
  \textbf{(a)}~SOC backbone: test images land near correct vehicle clusters.
  \textbf{(b)}~EOC backbone: test images scatter randomly, explaining the
  43.7\% baseline.}
  \label{fig:umap}
\end{figure}

Table~\ref{tab:ablation} shows the effect of progressively adding training
angles. Adding 60\textdegree{} alone accounts for +30.4~pp---the majority of
the total +42.4~pp gain from all additional angles. Large angular contrast
forces the model to learn angle-invariant representations.

\begin{table}[t]
  \caption{Angle Ablation: EOC Test Accuracy (\%)}
  \label{tab:ablation}
  \centering
  \begin{tabular}{lcc}
    \toprule
    Training Angles & Accuracy & $\Delta$ vs.\ Baseline \\
    \midrule
    15\textdegree{} only (EOC baseline) & 43.7 & --- \\
    15\textdegree{} + 60\textdegree{} & 74.1 & +30.4 \\
    15\textdegree{} + 30\textdegree{} + 60\textdegree{} & 78.3 & +34.6 \\
    All 4 angles ($N$=200) & 86.1 & +42.4 \\
    \bottomrule
  \end{tabular}
\end{table}

\subsection{Data Efficiency}

Table~\ref{tab:efficiency} shows accuracy versus $N$ images per (vehicle,
angle) group, tested on held-out 15\textdegree{} images. The key finding is
the \textit{crossover at $N$=25}: multi-angle training with only 25 images per
vehicle per angle already surpasses the EOC baseline (43.7\%) achieved by a
model trained on thousands of single-angle images. Angle diversity is more
valuable than data volume. The curve flattens above $N$=400 (0.4~pp gain from
400 to 500), suggesting diminishing returns beyond this threshold.

\begin{table}[t]
  \caption{Data Efficiency: Test Accuracy (\%) vs.\ $N$ per (Vehicle, Angle)}
  \label{tab:efficiency}
  \centering
  \begin{tabular}{ccc}
    \toprule
    $N$ & Accuracy & vs.\ EOC Baseline \\
    \midrule
    5   & 10.8 & --- \\
    10  & 21.7 & --- \\
    25  & 49.0 & +5.3 \\
    50  & 65.8 & +22.1 \\
    100 & 77.0 & +33.3 \\
    200 & 86.1 & +42.4 \\
    400 & 89.2 & +45.5 \\
    500 & 89.6 & +45.9 \\
    Full & 97.2 & +53.5 \\
    \bottomrule
  \end{tabular}
\end{table}

\subsection{Negative Results: What Does Not Work}

\textbf{Angle-weighted sampling.} Oversampling rarer angles using
inverse-frequency weights is consistently worse than uniform sampling
(Table~\ref{tab:negative}). Weighting does not help because the bottleneck is
data volume per angle, not sampling bias.

\textbf{Test-time augmentation (TTA).} TTA with $k$=7 augmented views reduces
accuracy from 89.8\% to 65.6\% ($-$24.3~pp). In optical vision, geometric
augmentations produce semantically equivalent views; in SAR, backscatter is
highly sensitive to azimuth and depression angle, making augmented views
out-of-distribution. Angle invariance cannot be approximated through
augmentation---it requires actual multi-angle training data.

\begin{table}[t]
  \caption{Negative Results: Weighted Sampling ($N$=200)}
  \label{tab:negative}
  \centering
  \begin{tabular}{cccc}
    \toprule
    $N$ & Baseline & Weighted & $\Delta$ \\
    \midrule
    25  & 40.8 & 40.8 & 0.0 \\
    50  & 55.2 & 52.5 & $-$2.7 \\
    100 & 70.4 & 68.1 & $-$2.3 \\
    200 & 86.1 & 80.1 & $-$6.0 \\
    \bottomrule
  \end{tabular}
\end{table}

\subsection{Few-Shot Prototypical Classification}

Table~\ref{tab:proto} shows prototypical classification accuracy using the
frozen SOC backbone. On ATRNet-STAR, 5 reference images per class achieves
98.2\%---within 0.3~pp of using 100 references. The curve saturates
immediately, demonstrating that a well-trained backbone produces features so
discriminative that nearest-centroid classification requires minimal
calibration data. DINOv2, by contrast, achieves near-chance performance
(6.1\% for 40 classes, random chance $=$ 2.5\%), confirming that optical
pretraining is orthogonal to SAR ATR. On MSTAR, a \textit{single} reference
image per class achieves 98.3\%, statistically indistinguishable from the
99.1\% SOC ceiling.

\begin{table}[t]
  \caption{Few-Shot Prototypical Accuracy (\%) vs.\ $N$ Reference Images/Class}
  \label{tab:proto}
  \centering
  \begin{tabular}{ccccc}
    \toprule
    $N$ & ATRNet-STAR & DINOv2 & MSTAR \\
    \midrule
    1  & ---  & ---  & 98.3 \\
    5  & 98.2 & 6.1  & 99.1 \\
    10 & 98.4 & 6.9  & 99.2 \\
    25 & 98.4 & 8.4  & 99.2 \\
    50 & 98.5 & 8.6  & 99.1 \\
    \bottomrule
  \end{tabular}
\end{table}

\subsection{Cross-Dataset Transfer}

Using the ATRNet-STAR backbone as a feature extractor on MSTAR without any
MSTAR training, prototypical classification reaches 61.1\% with full reference
sets versus 99.1\% for the MSTAR-specific backbone---a 38~pp gap quantifying
the SAR-to-SAR domain shift. Placed in context: DINOv2 (optical pretraining)
achieves $\sim$12\%; ATRNet-STAR (SAR, different targets) achieves 61.1\%;
MSTAR-specific training achieves 99.1\%. Transfer accuracy scales with domain
proximity: optical$\to$SAR transfer is negligible; SAR$\to$SAR transfer is
meaningful but insufficient for deployment.

% ── V. Discussion ─────────────────────────────────────────────────────────────
\section{Discussion}

\textbf{Why does 60\textdegree{} add so much?} The large angular contrast
between 60\textdegree{} and the 15\textdegree{} test condition forces the
network to discover angle-invariant features rather than angle-specific
artifacts, acting as a strong regularizer.

\textbf{Why does TTA fail in SAR?} In optical vision, geometric augmentations
produce semantically equivalent views. In SAR, backscatter intensity is
highly sensitive to azimuth and depression angle; rotated chips fall outside
the training distribution. This distinguishes SAR ATR from most optical tasks.

\textbf{Practical implications.} For new SAR ATR deployments: (1)~train a
backbone on whatever labeled data is available, even from a different dataset;
(2)~use prototypical classification for adaptation with a handful of reference
images; (3)~prioritize collecting reference images at multiple depression
angles over maximizing count at a single angle.

\textbf{Limitations.} All experiments use single-polarization magnitude
images. Dual-polarization and complex-valued inputs likely carry additional
discriminative information. Both datasets are collected in controlled
conditions; performance on operational SAR imagery with extended clutter
and non-cooperative targets may differ.

% ── VI. Conclusion ─────────────────────────────────────────────────────────────
\section{Conclusion}

We characterized SAR ATR generalization along three axes---depression angle
diversity, data quantity, and cross-dataset transfer---with six findings:
(1)~adding 60\textdegree{} to a 15\textdegree{}-only baseline yields +30.4~pp,
more than half the total gain from all additional angles;
(2)~25 images per vehicle per angle already surpasses thousands of single-angle
images, with diminishing returns above $N$=400;
(3)~weighted sampling and TTA both fail, with TTA actively hurting
($-$24.3~pp);
(4)~a frozen SOC backbone enables 98\%+ few-shot accuracy with 5 reference
images, replicated across ATRNet-STAR and MSTAR;
(5)~SAR-to-SAR transfer reaches 61.1\%, far above optical transfer but
below task-specific training; and
(6)~DINOv2 achieves $\sim$6--12\% on a 40-class SAR problem, motivating
SAR-specific pretraining.

% ── Acknowledgments ───────────────────────────────────────────────────────────
\section*{Acknowledgments}
Code implementation and manuscript preparation were assisted by Claude
(Anthropic). All experimental design, interpretation, and scientific claims
are the sole responsibility of the author.

% ── References ────────────────────────────────────────────────────────────────
\begin{thebibliography}{1}

\bibitem{atrnet}
Y.~Liu \textit{et al.}, ``ATRNet-STAR: A Large Dataset and Benchmark Towards
Remote Sensing Object Recognition in the Wild,'' submitted to \textit{IEEE
Trans. Pattern Anal. Mach. Intell.}

\bibitem{mstar}
V.~R.~Lalmani, ``MSTAR Dataset 10 Classes,'' IEEE DataPort, 2025,
doi:~10.21227/bfp9-7019.

\bibitem{convnext}
Z.~Liu \textit{et al.}, ``A ConvNet for the 2020s,'' in \textit{Proc.\
IEEE/CVF CVPR}, 2022, pp.~11976--11986.

\bibitem{dinov2}
M.~Oquab \textit{et al.}, ``DINOv2: Learning Robust Visual Features without
Supervision,'' \textit{arXiv:2304.07193}, 2023.

\bibitem{proto}
J.~Snell, K.~Swersky, and R.~S.~Zemel, ``Prototypical Networks for Few-shot
Learning,'' \textit{arXiv:1703.05175}, 2017.

\bibitem{dgpnet}
X.~Zhou, Q.~Wei, and Y.~Zhang, ``SAR Target Depression Angle Invariant
Recognition of Few-Shot Learning Via Dense Graph Prototype Network,''
\textit{IEEE J. Sel. Topics Appl. Earth Observ. Remote Sens.}, Jan.~2025,
doi:~10.1109/JSTARS.2025.3586129.

\bibitem{sarx}
W.~Li \textit{et al.}, ``SARATR-X: Towards Building A Foundation Model for
SAR Target Recognition,'' 2025.

\end{thebibliography}

\end{document}
